\chapter{Introduction}
\label{ch:intro}

\section{Motivation}

Many questions about tables cannot be answered from a single cell. ``What is the middle name of the player with the second most National Football League career rushing yards?'' requires finding the row whose rank is 2 in a table of rushing leaders, following the link from that row's player cell to the player's Wikipedia article, and reading the middle name there (Figure~\ref{fig:example}). Each step is a decision about which piece of evidence to look at next, and a wrong choice at the first step cannot be repaired by a good reader at the end. Benchmarks such as HybridQA \citep{chen-etal-2020-hybridqa}, OTT-QA \citep{chen2021open} and MMQA \citep{wu2025mmqa} were built around exactly this kind of chained evidence, across rows, linked passages and joined tables.

Existing table reasoning systems all contain such a selection step, but none of them learns it from the quality of the final answer. HYBRIDER, the reference model for HybridQA, ranks cells with a classifier trained on weak labels and keeps only the top cell at each of two greedy steps \citep{chen-etal-2020-hybridqa}. TableRAG retrieves columns and cells by embedding similarity \citep{chen2024tablerag}. Chain-of-Table and Dater ask a large language model to pick rows and columns greedily \citep{wang2024chainoftable,ye2023dater}. The hybrid-graph method ODYSSEY prunes its graph with a fixed similarity threshold and a breadth-first traversal of fixed depth \citep{agarwal-etal-2025-hybrid}. All of these choose the next piece of evidence by how relevant it looks now, not by how likely it makes a correct final answer, and all of them stop after a fixed number of steps. When the stages are evaluated separately the cost of this is visible: in HYBRIDER, stage accuracies of 87--89\% for linking, ranking and hopping compound before the reader is even reached.

The decision ``which entity next, or stop'' is a sequential decision with a delayed reward, which is the setting reinforcement learning was designed for. Reinforcement learning has been used for path finding in knowledge graphs \citep{xiong2017deeppath,das2018minerva,lin2018multihop}, for selecting a region inside a single table \citep{lee2024learningtoreduce,wu2025tabler1region}, and for training search agents over text \citep{jin2025searchr1,jiang2025s3,wang2025stepsearch}. The one attempt to apply it to table--text question answering chose between retrieval modules rather than entities and converged to a fixed strategy that did no better than always retrieving text \citep{jose2024qatexttablesdrl}. Whether a learned, entity-level selector can do better than the heuristics used by current frameworks is therefore an open question.

\section{Problem statement}

The project studies multi-hop table question answering in which a policy selects evidence one entity at a time. At every step the policy observes the question and the evidence gathered so far, and chooses either one entity from a candidate set (a table row or a passage linked from the table) or to stop; a frozen reader then answers from the selected evidence. The policy is trained with Proximal Policy Optimisation \citep{schulman2017ppo} in an actor--critic setup with generalised advantage estimation \citep{schulman2016gae}, after a behaviour-cloning warm start, on a reward that combines the reader's answer quality, a cost per selected entity and an evidence-recall signal. The selector is meant to be a drop-in replacement for the heuristic selection step of existing table reasoning pipelines, so that the reasoner itself stays unchanged.

\section{Objectives}

The work has four objectives:
\begin{enumerate}
  \item to survey the literature on learned evidence selection for multi-hop and table question answering, identify the most relevant papers, and state precisely what they leave open (Chapters~\ref{ch:literature} and~\ref{ch:gaps});
  \item to formulate entity selection over tables and linked text as a Markov decision process with a well-defined state, action set and reward, including the behaviour of the evidence reward under different optimisers (Chapter~\ref{ch:formulation});
  \item to design and implement the policy, the PPO training procedure and the comparison methods, together with a plan that addresses each identified gap (Chapters~\ref{ch:method} and~\ref{ch:implementation});
  \item to obtain preliminary results on HybridQA that test the central claims against strong heuristics, with enough seeds and statistical testing for the comparisons to be meaningful (Chapters~\ref{ch:setup} and~\ref{ch:results}).
\end{enumerate}
The research questions that the experiments address are listed at the end of Chapter~\ref{ch:gaps}.

\section{Contributions of this report}

The report makes four contributions.
\begin{enumerate}
  \item \textbf{A focused survey and gap analysis.} More than one hundred candidate papers were checked and the relevant ones organised into four groups; the twelve most relevant are ranked and analysed, and seven precise research gaps are derived. The sharpest gap is that no published RL method selects table entities step by step with a learned stop, and the only RL controller for table--text QA did no better than a fixed strategy.
  \item \textbf{A formulation with two properties of the evidence reward.} Entity selection over a per-question graph of rows and linked passages is defined as an MDP with a reader-in-the-loop reward. We show that the evidence-recall signal leaves every return unchanged when it enters as a potential, and that group-relative methods such as GRPO cannot see it at all in that form.
  \item \textbf{A tested implementation and a plan.} A pointer actor--critic with PPO and GAE, behaviour cloning on set-valued targets, REINFORCE, GRPO, heuristic baselines and an oracle are implemented with unit tests and a resumable pipeline. All training and evaluation runs of this report together take 1.2 hours on one T4 GPU.
  \item \textbf{Preliminary results with statistical tests.} On the full HybridQA dev set, with the same frozen reader for every selector, PPO trained from scratch reaches \fScratch{} F1, against \fBmThree{} for the best heuristic and \fBc{} for behaviour cloning (paired-bootstrap gains of \dScratchBm{} and \dScratchBc{} points, $p<0.001$). Warm-starting PPO from the cloned policy, which the literature recommends, removed the gain, and GRPO degraded. These findings change the plan for the rest of the project.
\end{enumerate}


\section{Organisation}

Chapter~\ref{ch:background} summarises the reinforcement learning methods used. Chapter~\ref{ch:literature} reviews the related work and Chapter~\ref{ch:gaps} ranks the most relevant papers and derives seven research gaps. Chapter~\ref{ch:formulation} defines the MDP and the reward, Chapter~\ref{ch:method} the policy and the training algorithms, and Chapter~\ref{ch:implementation} the implementation plan. Chapters~\ref{ch:setup} and~\ref{ch:results} describe the experimental setup and the preliminary results, Chapter~\ref{ch:discussion} discusses their limitations, and Chapter~\ref{ch:conclusion} concludes.
